\section{The source GCT coefficient on the retained arithmetic cover}
\label{sec:conic-arithmetic}
The GCT IV calculation quoted and checked in the preceding source
section uses the Laurent parameter $q$ and the original coefficient
\[
 C(q)=q^{10}-q^4-q^{-4}+q^{-10}
 =(q-q^{-1})^2[7]_q[3]_q,\qquad
 [j]_q=\frac{q^j-q^{-j}}{q-q^{-1}}.
\]
Here $[j]_q$ denotes its Laurent-polynomial expression at $q=\pm1$
as well; division in the displayed expression is an identity away
from these points, not deletion of those points from the polynomial.
We now construct an exact map to the original retained cover.

\begin{proposition}[A base change with all coordinates retained]
\label{prop:conic-base-change}
For a field $k$ of characteristic different from $2$, set
\[
 D=k[a,r]/(r^2+a),\qquad
 E=k[p]\otimes_{k[a]}D,\quad a\longmapsto4-p^2.
\]
There is an isomorphism
\[
 E=k[p,r]/(p^2-r^2-4)\ \simeq\ k[q,q^{-1}]
\]
with maps
\[
 p\longmapsto q+q^{-1},\quad r\longmapsto q-q^{-1},
 \qquad q\longmapsto(p+r)/2,\quad q^{-1}\longmapsto(p-r)/2.
\]
Its original base coordinate is exactly
$a=4-p^2=-(q-q^{-1})^2$.
\end{proposition}
\begin{proof}
Tensoring the displayed presentation of $D$ replaces $a$ by $4-p^2$,
giving the stated relation. The Laurent expressions satisfy that
relation. Conversely $(p+r)(p-r)/4=1$, so the proposed inverse images
of $q,q^{-1}$ have product one and define a Laurent ring map.
The sums and differences of these expressions recover $p,r$.
Starting with $q$, the same compositions recover $q,q^{-1}$.
Thus the maps are inverse on generators, proving the isomorphism
and every stated base-coordinate identity.
\end{proof}

The variable $p$ is new; it does not rename or replace the original
third source coordinate $w$. On the original simple-root chart with
$b=c=0$, the retained formulas become
\[
 \alpha=-2r,\quad x=-\frac1{2r},\quad y=3r,\quad w=26r^2,
 \qquad a=-r^2,\quad r\ne0.
\]
Indeed substitution into $\alpha=P'/2$ and the full inverse formulas
gives each equality with its original coefficient.
Substitution into the original $F$ then gives $(a,b,c)=(-r^2,0,0)$.
After Proposition~\ref{prop:conic-base-change}, these are rational
functions of $q$ defined precisely on $q\ne0,\pm1$:
$q=0$ was already absent from the Laurent parameter space, and
$q=\pm1$ is equivalent to $r=0$.
For each such $q$, the finite boundary point $(0,0,a)$ in the original
fibre is retained separately. It is not a point of this $x\ne0$ chart.

The exact positive Laurent factor has the polynomial expression
\[
 [7]_q=p^6-5p^4+6p^2-1,\qquad [3]_q=p^2-1.
\]
For a direct verification, set $U_0=1,U_1=p$ and
$U_j=pU_{j-1}-U_{j-2}$. Substituting $p=q+q^{-1}$ makes
$U_j=q^j+q^{j-2}+\cdots+q^{-j}$, by cancellation of the interior
terms in the recurrence. Applying this through $j=6$ proves the
two displayed polynomials.
Consequently the original coefficient, without discarding its
vanishing factor, is
\begin{align}
 C(q)&=-a\bigl((4-a)^3-5(4-a)^2+6(4-a)-1\bigr)(3-a)\notag\\
     &=-a(7-14a+7a^2-a^3)(3-a).
 \label{eq:C-a}
\end{align}
The second expression follows by expanding the first, which remains
displayed to record the substitution.
The coefficient $-a$ is part of the original object in both expressions.
At $q=1$, $a=-(q-1)^2(q+1)^2/q^2$.
The other factors take values $7$ and $3$, so
\[
 \left.\frac{C(q)}{(q-1)^2}\right|_{q=1}=4\cdot7\cdot3=84.
\]
In characteristic zero this proves a zero of order exactly two.
It proves a zero of order one in the original base coordinate $a$:
the coefficient of $a$ in \eqref{eq:C-a} is $-21$.
The order change is precisely the displayed ramification, not a
lost specialization layer.

\subsection{The resulting rank-four arithmetic extension}
Now take $k=\mathbb C$ and retain $a=2s$.
The same ring is finite free over $\mathbb C[s]$:
\begin{equation}\label{eq:quartic-q}
 \mathbb C[q,q^{-1}]
 \simeq\mathbb C[s,q]/\bigl(q^4+(2s-2)q^2+1\bigr).
\end{equation}
To prove this, the relation gives
$q^{-1}=-q^3+(2-2s)q$, so $q$ is already invertible in the quotient.
Conversely the Laurent map $s=-(q-q^{-1})^2/2$ gives the quartic.
The two substitutions recover $s$ and $q$, proving inverse ring maps.
Division by the monic quartic gives the free basis $(1,q,q^2,q^3)$.
In the conic presentation the alternative basis $(1,r,p,rp)$ is
free over $\mathbb C[s]$ by the two monic relations
$r^2=-2s$, $p^2=4-2s$.
The conversion identities are
\[
 r=q^3+(2s-1)q,\quad
 p=-q^3+(3-2s)q,\quad rp=2q^2+2s-2,\quad q=(p+r)/2.
\]
They follow from the expression for $q^{-1}$ and prove the basis
conversion in both directions (also $q^2=(rp+2-2s)/2$).

For the actual topological $\mathbb C[s]$-module $(Q,S)$ used in
the retained arithmetic construction, form
$Q^{(4)}=\mathbb C[q,q^{-1}]\otimes_{\mathbb C[s]}Q$.
Give it the fourfold direct-sum topology in the monic basis.
Multiplication by $q$ has the continuous block matrix
\begin{equation}\label{eq:quartic-operator}
 K=
 \begin{pmatrix}
 0&0&0&-I\\
 I&0&0&0\\
 0&I&0&2I-2S\\
 0&0&I&0
 \end{pmatrix},\qquad
 \widehat S=\operatorname{diag}(S,S,S,S).
\end{equation}
The last column is exactly the quartic relation. Thus
$K^4+(2\widehat S-2I)K^2+I=0$ and
$K^{-1}=-K^3+(2I-2\widehat S)K$.
Both are continuous polynomial expressions in commuting block
operators. The two original root and trace coordinates are
\[
 R=K-K^{-1},\quad P=K+K^{-1},\qquad
 R^2=-2\widehat S,\quad P^2=4I-2\widehat S.
\]
These identities follow either by the proved ring maps or by direct
substitution of the inverse polynomial.
In the alternative ordering $(1,r,p,rp)$, multiplication by $r$
is the direct sum of two copies of
$\left(\begin{smallmatrix}0&-2S\\I&0\end{smallmatrix}\right)$.
Thus this is the exact two-stage extension of the parent's two-sheet
module, with its factors and topology retained.
Projection to the constant coefficient and its section show, by the
same inverse construction used earlier, that
$\operatorname{Spec}\widehat S=\operatorname{Spec}S$.

Equation~\eqref{eq:C-a} also supplies a completely specified operator:
\begin{equation}\label{eq:coefficient-operator}
 C(K)=
 -2\widehat S(7I-28\widehat S+28\widehat S^2-8\widehat S^3)
 (3I-2\widehat S).
\end{equation}
Laurent evaluation is legitimate because $K$ has the explicit inverse.
This proves a concrete transport of a source GCT coefficient into the
retained arithmetic module. It does not assert that all nonstandard
quantum-group generators act on $Q^{(4)}$: only this scalar Laurent
parameter algebra and its proved action have been constructed.

\subsection{The original coefficient heat flow on this cover}
Use material time $\tau$, which generates the forward coefficient heat
flow of the retained polynomial, with $a'=2$, $b'=6c$, $c'=0$.
The arithmetic heat parameter $t$ acts instead by $-\partial_s^2/4$;
Section~\ref{sec:full-material-interface} proves its exact pullback and
retains the original Euclidean diffusion coefficients.
For the entire original cubic,
$\partial_\tau(cr^3-2r^2+br-2a)=6cr-4=\partial_r^2P$.
On $b=c=0$ this is $P=-2(r^2+a)$ with $a'=2$.
Implicit differentiation at a simple root gives $r'=-1/r$.
On the conic the second relation $p^2=4-a$ gives $p'=-1/p$
when $p\ne0$. Therefore the exact lifted coefficient derivation is
\[
 D_\tau=-\frac{q}{rp}\frac{d}{dq},\qquad
 D_\tau a=2,\quad D_\tau s=1,\quad
 D_\tau r=-\frac1r,\quad D_\tau p=-\frac1p.
\]
Indeed $dr/dq=1+q^{-2}=p/q$ and $dp/dq=1-q^{-2}=r/q$;
multiplication by the displayed coefficient proves all four identities.
Its domain is the Laurent curve with $rp\ne0$, equivalently
$q\notin\{0,1,-1,i,-i\}$ over $\mathbb C$.
The points $q=\pm1$ have the original excluded root $r=0$.
At $q=\pm i$ one has $p=0$ and $r=\pm2i$; these are the new
ramification points of the additional $p$-cover over $s=2$,
while the original $r$-chart remains finite.
Thus both exceptional loci and their different coordinate causes
are recorded by explicit equations.

For the original source coordinates this heat derivative is
\[
 D_\tau x=-\frac1{2r^3},\qquad
 D_\tau y=-\frac3r,\qquad
 D_\tau w=-52.
\]
Differentiating the proved identity gives
\[
 F(-1/(2r),3r,26r^2)=(-r^2,0,0),\qquad
 DF\,(D_\tau x,D_\tau y,D_\tau w)^{\mathsf T}=(2,0,0)^{\mathsf T}.
\]
Since $\det DF=-2$, this is exactly the original inverse-Jacobian
heat lift on the stated slice. The boundary point has derivative
$(0,0,2)$ separately.
For the GCT source-displayed coefficient polynomial the chain rule, with the full
polynomial $C_a(a)$ in \eqref{eq:C-a}, gives
$D_\tau C(q)=2C_a'(a)$. No analytic unit or vanishing factor is
deleted in this derivation. This is an action on the localized
coefficient ring. A differentiation action on the topological
arithmetic quotient is not inferred from its multiplication action;
the latter is exactly the operator action already proved in
\eqref{eq:quartic-operator}--\eqref{eq:coefficient-operator}.

The source identifies the classical spectral coordinate through
$\Xi(s)=\xi(1/2+is)=8H_0(2s)$.
Nothing in the finite-free extension changes that original coordinate.
A certified off-real spectral point of $S$ would remain off-real in
each direct summand; a nonreal eigenvalue of $K$ alone has no such
implication. The exact identity $\widehat S=-(K-K^{-1})^2/2$ gives
$S=\pi[-(K-K^{-1})^2/2]\iota$, with constant-coefficient projection
$\pi:Q^{(4)}\to Q$ and section $\iota:Q\to Q^{(4)}$.
This is the required test for the original coordinate.
None of these identities supplies
an off-real $s$, a positivity proof for all GCT coefficients, or a
complexity lower bound. They supply the exact common algebra and the
specific transported coefficient on which further calculations can act.

\subsection{The exact positivity test under two real structures}
On $\mathbb C[q,q^{-1}]$ define two conjugate-linear involutions by
$q^{*_{\rm real}}=q$ and $q^{*_{\rm circle}}=q^{-1}$, conjugating
complex coefficients in each case. Multiplicativity on Laurent
monomials and the fact that applying either map twice fixes $q$
prove both involution identities on the entire commutative algebra.
The conic isomorphism gives
\[
 p^{*_{\rm real}}=p,\quad r^{*_{\rm real}}=r,\qquad
 p^{*_{\rm circle}}=p,\quad r^{*_{\rm circle}}=-r.
\]
Both fix $a=-r^2$ and $s=a/2$. Thus they extend the same real structure
on the arithmetic base with different root-coordinate adjoints.

Evaluation at $q_0\ne0$ preserves $*_{\rm real}$ exactly when
$\overline{q_0}=q_0$, and preserves $*_{\rm circle}$ exactly when
$\overline{q_0}=q_0^{-1}$. These say, respectively, that $q_0$ is
nonzero real and that $|q_0|=1$.
For positive real $q_0$, all monomials of $[7]_q[3]_q$ are positive
and $(q-q^{-1})^2\ge0$, so $C(q_0)\ge0$.
For a precise circle character instead take
\[
 q_0=\frac{\sqrt{14}+i\sqrt2}{4},\quad
 q_0^{-1}=\frac{\sqrt{14}-i\sqrt2}{4},\quad
 p_0=\frac{\sqrt{14}}2,\quad r_0=\frac{i}{\sqrt2}.
\]
Products and sums verify $|q_0|=1$, $a_0=1/2$, $s_0=1/4$.
The full coefficient in \eqref{eq:C-a} then gives
\[
 C(q_0)=-\frac12
 \left(7-7+\frac74-\frac18\right)\frac52=-\frac{65}{32}.
\]
This exactly evaluates the source's mixed-sign coefficient, retaining
its positive residual Laurent factor and signed vanishing factor.
It rules out the particular claim of nonnegativity at every circle
character. It proves neither a failure of the source's stated
canonical-basis positivity nor a negative value on an established
zeta spectral point: $s_0=1/4$ has not been asserted to belong to that
spectrum. The maps and involutions above specify the entire comparison.
